\documentclass[11pt,a4paper]{article}
\usepackage[margin=0.85in,top=1in,bottom=1in]{geometry}
\usepackage{amsmath,amssymb,amsfonts}
\usepackage{graphicx}
\usepackage{float}
\usepackage{enumitem}
\usepackage{microtype}
\usepackage{caption}
\usepackage[hidelinks]{hyperref}
\usepackage[most,breakable]{tcolorbox}
\tcbuselibrary{breakable,skins}
\usepackage{fontspec}
\setmainfont{DejaVu Serif}
\setsansfont{DejaVu Sans}
\setmonofont{DejaVu Sans Mono}
\captionsetup{font=small,labelfont=bf,skip=4pt}
\setlist{leftmargin=1.5em,topsep=2pt}

%% Breakable problem card — prevents content cutoff across pages
\newtcolorbox{solcard}[3]{
  breakable, enhanced,
  colback=white, colframe=gray!50, arc=2pt,
  title={\textbf{#1}\hfill\textsf{\small #3\,pts}},
  fonttitle=\normalsize\sffamily\bfseries,
  before upper={\small\textit{#2}\par\smallskip\hrule height 0.4pt\smallskip},
  top=4pt, bottom=6pt, left=7pt, right=7pt,
  parbox=false,
}

\title{\textbf{Physics of raindrops --- Solution}}
\author{\normalsize X24 (2024) — English translation}
\date{}
\begin{document}\maketitle\vspace{-0.5em}
\begin{solcard}{A1}{Find the change in the free energy of water vapor if a drop of radius $r$ is formed from it. Express your answer in terms of $r$, $\sigma$, $\varphi$, $R$, $T$, $\rho_L$, $\mu$.}{1.00}

Due to the surface tension energy, the free energy will increase by $$ \Delta G_{surf} = \sigma A = 4 \pi \sigma r^2. $$On the other hand, for/at переходе vaporа в жидкое состояние его free energy уменьшается на $$ \Delta G _v = \nu R T \ln \varphi, $$where количество вещества в капле $$ \nu = \frac{4\pi \rho_L r^3}{3 \mu}. $$Here we have used the formula for changing the free energy of vapor and the fact that for saturated vapor the free energy is equal to the free energy of the liquid.

\par\smallskip\noindent\textbf{Answer:} \textbackslash{}textbf{Answer:} $$ \Delta G = 4 \pi \sigma r^2 - \frac{4\pi \rho_L }{3 \mu}r^3 R T \ln \varphi $$
\end{solcard}

\begin{solcard}{A2}{Find the critical value of the droplet radius $r_c$ at which $\Delta G$ is maximum, as well as the corresponding value $\Delta G_с$. Express your answer using $\sigma$, $\varphi$, $R$, $T$, $\rho_L$, $\mu$. Find the numerical value of $r_c$ at $\varphi = 1.01$.}{0.80}

To find the maximum, we find the derivative $$ \frac{\partial \Delta G}{\partial r}  = 8 \pi \sigma r - \frac{4\pi \rho_L }{ \mu} r^2 RT \ln \varphi = 0. $$Hence $$ r_c = \frac{2 \sigma \mu}{\rho_L R T \ln \varphi}. $$Substituting в формулу for $\Delta G$, we get $$ \Delta G_c = \frac{16 \pi}{3} \frac{\sigma^3 \mu^2}{\rho_L^2 R^2 T^2 \ln^2 \varphi} = \frac{4\pi r_c^2 \sigma}{3}. $$

\par\smallskip\noindent\textbf{Answer:} \textbackslash{}textbf{Answer:} $$ r_c = \frac{2 \sigma \mu}{\rho_L R T \ln \varphi} = 1.15 \cdot 10^{-7}~\text{m}, \quad \Delta G_c = \frac{16 \pi}{3} \frac{\sigma^3 \mu^2}{\rho_L^2 R^2 T^2 \ln^2 \varphi}. $$
\end{solcard}

\begin{solcard}{A3}{Consider a drop of critical radius $r_c$. Determine the time $\tau$ during which the number of molecules in it will increase by $g$. Express your answer using $r_c$, $g$, $p_s$, $m$, $k$, $T$, $\varphi$. Assume that during the growth process the radius of the drop does not change; the evaporation of molecules from the drop can be neglected. It is known that $$ dN = dt dS  \frac{p_v}{\sqrt{2\pi m k T}} $$ молекул. Here $p_v$ – pressure vaporа, $m$ – mass молекул, $T$, the gas temperature, falls on the surface area $dS$ in time $dt$.}{0.70}

$$ dN = 4\pi r_c^2 \frac{p_v}{\sqrt{2\pi m kT}} = 4\pi r_c^2 \frac{p_s \varphi}{\sqrt{2\pi m kT}} $$молекул. Here использовано соratio $p_v = p_s \varphi$. Since изменением radiusа можно пренебречь, coefficient пропорциональности постоянен, and therefore искомое time $$ \tau = g \left( 4\pi r_c^2 \frac{p_s \varphi}{\sqrt{2\pi m kT}} \right)^{-1}. $$ hits the entire surface area of ​​the drop in time $dt$

\par\smallskip\noindent\textbf{Answer:} \textbackslash{}textbf{Answer:} $$ \tau = \frac{g \sqrt{2\pi m kT}}{4\pi r_c^2 p_s \varphi}. $$
\end{solcard}

\begin{solcard}{A4}{Find the number of drops $J$ that are formed per unit time in a unit volume of supersaturated water vapor. Express your answer in terms of $\sigma$, $\varphi$, $p_s$, $r_c$, $T$, $g$.}{0.60}

By condition, in time $\tau$ all nuclei in the volume turn into drops, so $$ J = \frac{n_c}{\tau} = \frac{4\pi r_c^2 p_s \varphi}{g \sqrt{2\pi m kT}} n\exp\left( -  \frac{16 \pi}{3 kT} \frac{\sigma^3 \mu^2}{\rho_L^2 R^2 T^2 \ln^2 \varphi}\right). $$Let us express концентрацию vaporа via/through pressure: $$ n = \frac{p_v}{k T} = \frac{p_s \varphi}{k T}, $$ we get

\par\smallskip\noindent\textbf{Answer:} \textbackslash{}textbf{Answer:} $$ J =  \frac{4\pi r_c^2 }{ \sqrt{2\pi m kT}} \frac{p_s^2 \varphi^2}{k T} \frac{1}{g}\exp\left( -  \frac{16 \pi}{3 kT} \frac{\sigma^3 \mu^2}{\rho_L^2 R^2 T^2 \ln^2 \varphi}\right) = \frac{4\pi r_c^2 }{ \sqrt{2\pi m kT}} \frac{p_s^2 \varphi^2}{k T} \frac{1}{g}\exp\left( -  \frac{4\pi r_c^2 \sigma}{3 k T}\right). $$
\end{solcard}

\begin{solcard}{A5}{From the results of the previous paragraph it follows that the rate of droplet formation very much depends on the steam supersaturation coefficient. Determine numerically the value of the vapor supersaturation coefficient $\varphi$, at which one drop per second is born at a temperature $T = 283~\text{K}$ in $1~\text{cm}^3$ air. Consider that $g = 100$. The remaining numerical data are given at the beginning of the problem.}{0.90}

Let's substitute the expression for $r_c$ into the result of the previous paragraph: $$ J = \frac{4\pi p_s^2}{\sqrt{2\pi m k T}} \frac{4 \sigma^2 \mu^2}{g k T \rho_L^2 R^2 T^2} \frac{\varphi^2}{\ln^2 \varphi}\exp\left( -  \frac{16 \pi}{3 kT} \frac{\sigma^3 \mu^2}{\rho_L^2 R^2 T^2 \ln^2 \varphi}\right) = J_0 \frac{\varphi^2}{\ln^2 \varphi} \exp\left(- \frac{A}{\ln^2 \varphi} \right), $$where $$ J_0 = \frac{16\pi p_s^2}{\sqrt{2\pi m k T}} \frac{ \sigma^2 \mu^2}{g k T \rho_L^2 R^2 T^2} =  2.37 \cdot 10^{30} ~ \text{m}^{-3} \cdot \text{s}^{-1}, $$$$ A = \frac{16 \pi}{3 kT} \frac{\sigma^3 \mu^2}{\rho_L^2 R^2 T^2 } =  106. $$Нам нужно поrayandть value $J = 10^6 ~ \text{m}^{-3} \cdot \text{s}^{-1}$. Численно we find $\varphi \approx 3.86 $. For/Atведем also таблицу значений J for/at близких значениях $\varphi$, we see that growth is happening very quickly.


$\varphi$ $J, \, \text{m}^{-3} \cdot \text{s}^{-1}$ 3.5 $8.4 \cdot 10^1$ 3.6 $1.61 \cdot 10^2$ 3.7 $2.34 \cdot 10^4$ 3.8 $2.79 \cdot 10^5$ 3.9 $2.68 \cdot 10^6$

\par\smallskip\noindent\textbf{Answer:} \textbackslash{}textbf{Answer:} $$ \varphi = 3.86 $$
\end{solcard}

\begin{solcard}{B1}{For saturated vapor in equilibrium with a liquid, express the derivative of pressure with respect to temperature $dp_s/dT$ in terms of $p_s$, $L$, $R$, $T$, $\mu$. Using the result obtained, find the relative change in the density of saturated water vapor $\Delta \rho_s/\rho_s$ with a small change in temperature $\Delta T$. Express your answer using $\Delta T$, $T$, $L$, $\mu$, $R$. You can use the relationship between small changes in pressure, density and temperature of an ideal gas $$ \frac{\Delta p_s}{p_s} = \frac{\Delta \rho_s}{\rho_s} +\frac {\Delta T}{T}. $$}{0.80}

The dependence of saturated steam pressure on temperature is determined by the Clapeyron-Clausius equation (we assume that the volume of steam is much larger than the corresponding volume of water at the same temperature): $$ \frac{dp_s}{dT} = \frac{L \mu p}{R T^2}. $$Using соratio from the condition, let us find $$ \frac{\Delta \rho_s}{\rho_s} = \frac{\Delta p_s}{p_s} - \frac{\Delta T}{T} = \frac{\lambda \mu p}{R T^2} \Delta T  - \frac{\Delta T}{T} = \frac{\Delta T}{T} \left( \frac{\mu L}{R T} - 1\right). $$

\par\smallskip\noindent\textbf{Answer:} \textbackslash{}textbf{Answer:} $$ \frac{\Delta \rho_s}{\rho_s} = \frac{\Delta T}{T} \left( \frac{\mu L}{R T} - 1\right). $$
\end{solcard}

\begin{solcard}{B2}{Express $dQ/dt$ in terms of $dM/dt$ and $L$.}{0.20}

Since heat is released only due to condensation of water $dQ = L dM$

\par\smallskip\noindent\textbf{Answer:} \textbackslash{}textbf{Answer:} $$ \frac{dQ}{dt} = L \frac{dM}{dt} $$
\end{solcard}

\begin{solcard}{B3}{Using the result of the previous paragraph and the heat equation, express the temperature difference between the drop and the atmosphere, $T_r- T$, in terms of $dM/dt$, as well as $r$, $L$, $K$.}{0.30}

From the heat equation $$ T_r -  T = \frac{1}{4 \pi r K} \frac{dQ}{dt} = \frac{L}{4 \pi r K} \frac{dM}{dt}. $$

\par\smallskip\noindent\textbf{Answer:} \textbackslash{}textbf{Answer:} $$ T_r -  T = \frac{L}{4 \pi r K} \frac{dM}{dt}. $$
\end{solcard}

\begin{solcard}{B4}{We assume that near the surface of the drop the density of water vapor is equal to the density of saturated vapor at the temperature of the drop. Assuming the temperature and density differences are small and using the results $B1$, $B3$, express the ratio $(\rho_r - \rho_s)/\rho_s$ ($\rho_r$ is the vapor pressure near the surface of the drop) in terms of $L$, $r$, $K$, $\mu$, $R$, $T$ and $dM/dt$.}{0.30}

From the result B1 $$ \frac{\rho_r - \rho_s}{\rho_s} = \frac{T_r - T}{T}  \left( \frac{\mu L}{R T} - 1\right). $$ Substituting the expression for the temperature difference into it, we obtain

\par\smallskip\noindent\textbf{Answer:} \textbackslash{}textbf{Answer:} $$ \frac{\rho_r - \rho_s}{\rho_s} =  \left( \frac{\mu L}{R T} - 1\right)\frac{L}{4 \pi r K T} \frac{dM}{dt}. $$
\end{solcard}

\begin{solcard}{B5}{Using the diffusion equation, express the ratio $(\rho_r - \rho_v)/\rho_s$ in terms of $dM/dt$, $r$, $D$, $\rho_s$.}{0.30}

From the diffusion equation $$ \rho_v - \rho_r = \frac{1}{4 \pi r D} \frac{dM}{dt}, $$

\par\smallskip\noindent\textbf{Answer:} \textbackslash{}textbf{Answer:} $$ \frac{\rho_r - \rho_v}{\rho_s} = - \frac{1}{4 \pi r \rho_s D} \frac{dM}{dt} $$
\end{solcard}

\begin{solcard}{B6}{By excluding the vapor density near the surface of the drop $\rho_r$ from the answers in the two previous paragraphs, you obtain an expression for $dM/dt$. Express your answer in terms of $\varphi$, $\mu$, $R$, $T$, $D$, $p_s$, $L$, $K$, $r$.}{0.60}

Subtracting the expressions from the last two points from each other, we get $$ \frac{\rho_v - \rho_s}{\rho_s} = \varphi - 1= \frac{1}{4 \pi r}\left( \left( \frac{\mu L}{R T} - 1\right)\frac{L}{ K T} +\frac{1}{\rho_s D}  \right) \frac{dM}{dt}. $$Also let us express density насыщенного vaporа via/through pressure: $$ \rho_s = \frac{\mu p_s}{R T}, $$we finally get

\par\smallskip\noindent\textbf{Answer:} \textbackslash{}textbf{Answer:} $$ \frac{dM}{dt} = \frac{4 \pi r (\varphi - 1)}{\left( \frac{\mu L}{R T} - 1\right)\frac{L}{ K T} +\frac{R T}{\mu p_s D} } $$
\end{solcard}

\begin{solcard}{B7}{The rate of increase in the radius of the drop has the form $$ \frac{dr}{dt} = \frac{\xi}{r^k}. $$ Determine $k$ and $\xi$, выtimesandте answer via/through $\varphi $, $\rho_L$, $\mu$, $R$, $T$, $D$, $p_s$, $L$, $K$.}{0.50}

The mass of the drop is related to the radius by the relation $$ M = \frac{4 \pi}{3} \rho_L r^3, $$therefore $$ \frac{dM}{dt} = 4 \pi \rho_L r^2 \frac{dr}{dt}, $$and therefore $$ \frac{dr}{dt} = \frac{1}{4 \pi \rho_L r^2} \frac{dM}{dt} = \frac{1}{r \rho_L}  \frac{\varphi - 1}{\left( \frac{\mu L}{R T} - 1\right)\frac{L}{ K T} +\frac{R T}{\mu p_s D} } $$

\par\smallskip\noindent\textbf{Answer:} \textbackslash{}textbf{Answer:} $$ k = 1, \quad \xi = \frac{\varphi - 1}{\left( \frac{\mu L}{R T} - 1\right)\frac{L}{ K T} +\frac{R T}{\mu p_s D} } \frac{1}{\rho_L}. $$
\end{solcard}

\begin{solcard}{B8}{Find the dependence of the radius of the drop on time. The initial radius of the drop is $r_0$. Express your answer in terms of $r_0$, $\xi$, $t$.}{0.50}

Let's integrate the equation $$ r \frac{dr}{dt} = \xi, $$we get $$ \frac{r^2}{2} - \frac{r_0^2}{2} = \xi t, $$

\par\smallskip\noindent\textbf{Answer:} \textbackslash{}textbf{Answer:} $$ r(t) = \sqrt{r_0^2 + 2 \xi t}. $$
\end{solcard}

\begin{solcard}{B9}{Let the initial radius of the drop be $r_0 = 0.7~\text{\mu{}m}$. Find the numerical value of the time during which it will grow to size $r_1 = 10~\text{\mu{}m}$ with a supersaturation coefficient $\varphi = 1.1$. The remaining numerical values ​​are given at the beginning of this part.}{0.50}

For the parameters from the condition $$ \xi = 9.04 \cdot 10^{-12} ~\text{m}^2 \cdot \text{s}^{-1}, $$then the time

\par\smallskip\noindent\textbf{Answer:} \textbackslash{}textbf{Answer:} $$ t = \frac{r_1^2 - r_0^2}{2 \xi} = 5.50~\text{s}. $$
\end{solcard}

\end{document}